% संपूर्ण मराठी ग्रंथातील प्रकरण 1: संच
% हा संपादनयोग्य स्रोतखंड आहे. संकलनासाठी संपूर्ण स्रोतसंग्रहातील
% अक्षररूपे, पूर्वभाग, चिन्हव्याख्या आणि आकृती-स्रोत आवश्यक आहेत.
% मूळ: Open Logic Project; मजकूर आणि रूपांतर CC BY 4.0.
% यंत्रानुवाद, दुरुस्ती आणि तपासणी: OpenAI Codex —
% GPT-5.6 Sol आणि GPT-6 Sol, दोन्ही Ultra effort.
% दुरुस्ती-पुनरावलोकन: GPT-6.1 Sol, Ultra effort.
\chapter{संच}








\section{विस्तारात्मकता}\label{sfr:set:bas:sec}

वस्तूंचा समूह एकच वस्तू {म्हणून} विचारात घेतला, की त्याला \emph{संच} म्हणतात.
संच बनवणाऱ्या वस्तूंना त्या संचाचे \emph{घटक} किंवा \emph{सदस्य} म्हणतात.
$x$ हा $A$ या संचाचा घटक असेल, तर आपण $x \in A$ असे लिहितो;
नसेल, तर $x \notin A$ असे लिहितो. ज्या संचात एकही घटक नसतो,
त्याला \emph{रिक्त} संच म्हणतात आणि तो ``$\emptyset$'' या चिन्हाने दर्शवतात.

\begin{explain}
आपण संच \emph{कसा निर्दिष्ट करतो}, त्यातील घटक \emph{कोणत्या क्रमाने}
मांडतो किंवा तेच घटक \emph{किती वेळा} मोजतो, याला महत्त्व नसते.
त्यात कोणते घटक आहेत, एवढेच महत्त्वाचे असते. ही गोष्ट आपण पुढील
तत्त्वात नेमकेपणाने मांडतो.
\end{explain}

\begin{defn}[विस्तारात्मकता]
  $A$ आणि $B$ हे संच असतील, तर $A = B$ तेव्हा आणि केवळ तेव्हाच, जेव्हा
  $A$ मधील प्रत्येक घटक हा $B$ चाही घटक असतो, आणि उलटही.
\end{defn}

विस्तारात्मकतेमुळे काही चिन्हपद्धती वापरणे योग्य ठरते. सर्वसाधारणपणे,
$a_{1}$, \dots, $a_{n}$ या वस्तू दिलेल्या असतील, तर
$\{a_{1}, \dots, a_{n}\}$ म्हणजे ज्यातील घटक हे
$a_1, \ldots, a_n$ आहेत \emph{तो विशिष्ट} संच होय. येथे ``\emph{तो विशिष्ट}''
या शब्दांवर भर आहे, कारण विस्तारात्मकतेनुसार असा \emph{एकच} संच असू शकतो.
विस्तारात्मकतेमुळे पुढील समानताही योग्य ठरते:
  \[    \{a, a, b\} = \{a, b\} = \{b,a\}.
  \]
म्हणजेच, संचांचा विचार करताना त्यातील घटक कोणत्या क्रमाने किंवा
किती वेळा निर्दिष्ट केले आहेत, याला महत्त्व नसते.


\begin{ex}
काही वस्तू दिलेल्या असतील, तर त्या एकत्र करून त्यांचा संच बनवता येतो.
उदाहरणार्थ, रिचर्डच्या भावंडांच्या संचात एक व्यक्ती आहे, आणि तो
$S=\{\textrm{Ruth}\}$ असा लिहिता येईल.
$4$ पेक्षा लहान धन पूर्णांकांचा संच $\{1, 2, 3\}$ आहे; पण तो
$\{3, 2, 1\}$ किंवा अगदी $\{1, 2, 1, 2, 3\}$ असाही लिहिता येतो.
विस्तारात्मकतेनुसार हे सर्व एकच संच आहेत. कारण $\{1, 2, 3\}$ मधील
प्रत्येक घटक हा $\{3, 2, 1\}$ चा (आणि $\{1,
2, 1, 2, 3\}$ चाही) घटक आहे, आणि उलटही.
\end{ex}


अनेकदा संचातील सर्व घटक यांना समान असलेला एखादा गुणधर्म सांगून आपण
संच निर्दिष्ट करू. त्यासाठी पुढील संक्षिप्त चिन्हपद्धती वापरू:
$\Setabs{x}{\phi(x)}$. येथे $\phi(x)$ हा असा गुणधर्म दर्शवतो की,
$x$ ची संचातील घटक मध्ये गणना होण्यासाठी त्यात तो गुणधर्म असला पाहिजे.


\begin{ex}
आपल्या उदाहरणातील $S$ हा संच पुढीलप्रमाणेही निर्दिष्ट करता आला असता:
\[S = \Setabs{x}{x \text{ हे रिचर्डचे भावंड आहे}}.
\]
\end{ex}



\begin{ex}
एखादी संख्या तिच्या स्वतःखेरीज इतर सर्व धन विभाजकांच्या बेरजेइतकी असेल तेव्हा
आणि केवळ तेव्हाच तिला \emph{परिपूर्ण} म्हणतात (हे विभाजक म्हणजे त्या संख्येला
निःशेष भाग जाणाऱ्या, पण तिच्याशी समान नसलेल्या संख्या).
उदाहरणार्थ, $6$ ही परिपूर्ण संख्या आहे, कारण तिचे असे विभाजक
$1$, $2$ आणि~$3$ आहेत, आणि $6 = 1 + 2 + 3$.
खरे तर, $10$ पेक्षा लहान धन पूर्णांकांपैकी $6$ हीच एकमेव परिपूर्ण संख्या आहे.
म्हणून विस्तारात्मकतेचा उपयोग करून आपण असे म्हणू शकतो:
\[	\{6\} = \Setabs{x}{x\text{ ही परिपूर्ण संख्या आहे आणि }0 \leq x \leq 10}
\]
उजवीकडील चिन्हलेखनाचे वाचन असे करतो: ``$x$ ही परिपूर्ण संख्या आहे आणि
$0 \leq x \leq 10$ अशा सर्व $x$ चा संच.'' संच कसा निर्दिष्ट केला आहे,
याला संचाचा विचार करताना महत्त्व नसते, ही गोष्ट वरील समानता स्पष्ट करते.
अधिक सर्वसाधारणपणे, $\phi(x)$ असणाऱ्या सर्व $x$ चा संच नेहमी एकच असतो,
याची विस्तारात्मकता खात्री देते.
त्यामुळे $\Setabs{x}{\phi(x)}$ याला $\phi(x)$ असणाऱ्या सर्व $x$ चा
\emph{तो विशिष्ट} संच म्हणणे विस्तारात्मकतेमुळे योग्य ठरते.
\end{ex}


संच समान असल्याचे दाखवण्याची पद्धत विस्तारात्मकतेमुळे मिळते:
$A = B$ हे दाखवण्यासाठी, $x \in A$ असेल तेव्हा $x \in B$ देखील असते,
आणि $y \in B$ असेल तेव्हा $y \in A$ देखील असते, हे दाखवा.

\begin{prob}
जास्तीत जास्त एकच रिक्त संच असतो हे सिद्ध करा. म्हणजे, $A$ आणि $B$ या
संचांत एकही घटक नसेल, तर $A = B$ हे दाखवा.
\end{prob}











\section{उपसंच आणि घातसंच}\label{sfr:set:sub:sec}

\begin{explain}
आपल्याला अनेकदा संचांची तुलना करावी लागेल. तुलना करण्याचा एक सरळ प्रकार
असा आहे: \emph{एका संचातील प्रत्येक वस्तू दुसऱ्या संचातही असते}.
ही स्थिती पुरेशी महत्त्वाची असल्यामुळे तिच्यासाठी आपण नवीन चिन्हपद्धतीची
ओळख करून घेऊ.
\end{explain}

\begin{defn}[उपसंच]
$A$ या संचातील प्रत्येक घटक हा $B$ चाही घटक असेल,
तर $A$ हा $B$ चा \emph{उपसंच} आहे असे म्हणतो, आणि $A \subseteq B$ असे लिहितो.
$A$ हा $B$ चा उपसंच नसेल, तर $A \not\subseteq B$ असे लिहितो.
$A \subseteq B$ पण $A \neq B$ असेल, तर $A \subsetneq B$ असे लिहितो
आणि $A$ हा $B$ चा \emph{उचित उपसंच} आहे असे म्हणतो.
\end{defn}

\begin{ex}
प्रत्येक संच स्वतःचाच उपसंच असतो, आणि $\emptyset$ हा प्रत्येक संचाचा उपसंच असतो.
सम नैसर्गिक संख्यांचा संच हा नैसर्गिक संख्यांच्या संचाचा उपसंच आहे.
तसेच, $\{ a, b \} \subseteq \{ a, b, c \}$.
पण $\{ a, b, e
\}$ हा $\{ a, b, c \}$ चा उपसंच नाही.
\end{ex}

\begin{ex}
$2$ ही संख्या पूर्णांकांच्या संचाचा घटक आहे, तर सम संख्यांचा संच
हा पूर्णांकांच्या संचाचा उपसंच आहे. मात्र, एखादा संच दुसऱ्या संचाचा
घटक आणि उपसंच \emph{दोन्हीही} असू शकतो. उदाहरणार्थ,
$\{0\} \in \{0, \{0\}\}$ आणि $\{0\} \subseteq \{0,
\{0\}\}$ देखील.
\end{ex}

विस्तारात्मकतेमुळे संचांच्या समानतेचा निकष मिळतो: $A = B$ तेव्हा आणि केवळ तेव्हाच,
जेव्हा $A$ मधील प्रत्येक घटक हा $B$ चाही घटक असतो, आणि उलटही.
``उपसंच'' या व्याख्येत $A \subseteq B$ याचा अर्थ या निकषाचा नेमका पहिला
भाग आहे: $A$ मधील प्रत्येक घटक हा $B$ चा घटक असतो.
अर्थात, $A$ आणि $B$ यांची अदलाबदल केली तरी ही व्याख्या लागू होते:
$B \subseteq A$ तेव्हा आणि केवळ तेव्हाच, जेव्हा $B$ मधील प्रत्येक घटक
हा $A$ चा घटक असतो. हाच विस्तारात्मकतेतील ``आणि उलटही'' हा भाग होय.
दुसऱ्या शब्दांत, संच एकमेकांचे उपसंच असतील तेव्हा आणि केवळ तेव्हाच ते समान
असतात, हे विस्तारात्मकतेवरून निष्पन्न होते.

\begin{prop}
$A = B$ तेव्हा आणि केवळ तेव्हाच, जेव्हा $A \subseteq B$ आणि $B \subseteq A$ दोन्ही असतात.
\end{prop}

आणखी काही उपयुक्त चिन्हपद्धतींची ओळख करून घेण्याची ही योग्य वेळ आहे.
$A$ हा $B$ चा उपसंच केव्हा असतो याची व्याख्या करताना आपण
``$A$ मधील प्रत्येक घटक हा \dots'' असे म्हटले, आणि त्या
``$\dots$'' च्या जागी ``$B$ चा घटक असतो'' असे घातले.
परंतु असा \emph{आकार} असणाऱ्या अभिव्यक्ती इतक्या नेहमी वापरल्या जातात की,
त्यांच्यासाठी आकारिक चिन्हपद्धती उपयोगी पडेल.

\begin{defn}\label{sfr:set:sub:forallxina}
$(\forall x \in A)\phi$ हे $\forall x(x \in A \lif
\phi)$ याचे संक्षिप्त रूप आहे. त्याचप्रमाणे, $(\exists x \in A)\phi$ हे
$\exists x(x
\in A \land \phi)$ याचे संक्षिप्त रूप आहे.
\end{defn}

या चिन्हपद्धतीचा उपयोग करून, $A \subseteq B$ तेव्हा आणि केवळ तेव्हाच, जेव्हा
$(\forall
x \in A)x \in B$, असे म्हणता येते.

आता आपण एका विशिष्ट प्रकारच्या संचाचा विचार करू: दिलेल्या संचाच्या सर्व
उपसंचांचा संच.

\begin{defn}[घातसंच]
$A$ या संचाच्या सर्व उपसंचांचा मिळून जो संच बनतो, त्याला
$A$ चा \emph{घातसंच} म्हणतात आणि तो $\Pow{A}$ असा लिहितात.
  \[    \Pow{A} = \Setabs{B}{B \subseteq A}
  \]
\end{defn}

\begin{ex}
$\{ a, b, c \}$ चे सर्व शक्य उपसंच कोणते? ते असे आहेत:
$\emptyset$, $\{a \}$, $\{b\}$, $\{c\}$, $\{a, b\}$, $\{a, c\}$, $\{b,
c\}$, $\{a, b, c\}$. या सर्व उपसंचांचा संच म्हणजे
$\Pow{\{a,b,c\}}$:
\[\Pow{\{ a, b, c \}} = \{\emptyset, \{a \}, \{b\}, \{c\}, \{a, b\},
\{b, c\}, \{a, c\}, \{a, b, c\}\}
\]
\end{ex}

\begin{prob}
$\{a, b, c, d\}$ चे सर्व उपसंच लिहा.
\end{prob}

\begin{prob}
$A$ मध्ये $n$ घटक असतील, तर $\Pow{A}$ मध्ये $2^n$
घटक असतात हे दाखवा.
\end{prob}











\section{काही महत्त्वाचे संच}\label{sfr:set:imp:sec}

\begin{ex}
आपण प्रामुख्याने ज्यातील घटक या गणिती वस्तू आहेत अशा संचांचा
विचार करणार आहोत. असे चार संच इतके महत्त्वाचे आहेत की त्यांना
विशिष्ट नावे दिलेली आहेत:
\begin{multline*}
    \Nat = \{0, 1, 2, 3, \ldots\} \\
    \shoveright{\text{नैसर्गिक संख्यांचा संच}}\\
    \shoveleft{\Int = \{\ldots, -2, -1, 0, 1, 2, \ldots\}} \\
    \shoveright{\text{पूर्णांकांचा संच}}\\
    \shoveleft{\Rat = \Setabs{\nicefrac{m}{n}}{m, n \in \Int\text{ आणि }n \neq 0}}\\
    \shoveright{\text{परिमेय संख्यांचा संच}}\\
    \shoveleft{\Real = (-\infty, \infty)}\\
    \text{वास्तव संख्यांचा संच (सांतत्य)}
\end{multline*}
हे सर्व \emph{अनंत} संच आहेत; म्हणजेच, प्रत्येकात अनंत घटक आहेत.

या संचांतून पुढे जाताना आपण आपल्या संग्रहात \emph{आणखी} संख्या समाविष्ट
करत असतो. खरे तर, $\Nat \subseteq \Int
\subseteq \Rat \subseteq \Real$ हे स्पष्ट असावे: कारण प्रत्येक नैसर्गिक संख्या
पूर्णांक असते; प्रत्येक पूर्णांक परिमेय असतो; आणि प्रत्येक परिमेय संख्या वास्तव असते.
त्याचप्रमाणे, $\Nat \subsetneq \Int \subsetneq
\Rat$ हेही स्पष्ट असावे, कारण $-1$ हा पूर्णांक आहे पण नैसर्गिक संख्या नाही,
आणि $\nicefrac{1}{2}$ ही परिमेय संख्या आहे पण पूर्णांक नाही.
$\Rat \subsetneq \Real$, म्हणजे परिमेय नसलेल्या काही वास्तव संख्या आहेत,
ही गोष्ट तितकी सहज स्पष्ट होत नाही.

कधीकधी आपण धन पूर्णांकांचा संच $\PosInt = \{1,
2, 3, \dots\}$ आणि फक्त पहिल्या दोन नैसर्गिक संख्या असणारा संच
$\Bin = \{0, 1\}$ यांचाही उपयोग करू.
\end{ex}


\begin{ex}[चिन्हमाला]
आणखी एक लक्षवेधी उदाहरण म्हणजे $A$ या वर्णमालेवरील \emph{सांत चिन्हमालांचा}
संच $A^{*}$: $A$ मधील घटकांचा कोणताही सांत अनुक्रम हा $A$ वरील चिन्हमाला
असतो. प्रत्येक $A$ वर्णमालेसाठी, $A$ वरील चिन्हमालांमध्ये आपण
\emph{रिक्त चिन्हमाला $\Lambda$} हिचाही समावेश करतो. उदाहरणार्थ,
\begin{multline*}
\Bin^*
=\{\Lambda,0,1,00,01,10,11,\\
000,001,010,011,100,101,110,111,0000,\ldots\}.
\end{multline*}
$x=x_{1}\ldots x_{n}\in A^{*}$ ही $A$ मधील $n$ ``अक्षरांनी'' बनलेली
चिन्हमाला असेल, तर तिची \emph{लांबी} $n$ आहे असे म्हणतो आणि $\len{x}=n$ असे लिहितो.
\end{ex}


\begin{ex}[अनंत अनुक्रम]
कोणत्याही $A$ संचासाठी आपण $A$ मधील घटक यांच्या अनंत अनुक्रमांचा
संच $A^\omega$ याचाही विचार करू शकतो. अनंत अनुक्रम
$a_1a_2a_3a_4\dots$ म्हणजे वस्तूंची एका दिशेने अनंत लांब जाणारी यादी;
यादीतील प्रत्येक वस्तू $A$ चा घटक असते.
\end{ex}











\section{संयोग आणि छेद}\label{sfr:set:uni:sec}

\begin{explain}
\ref{sfr:set:bas:sec} मध्ये आपण गुणधर्मानुसार संचांची व्याख्या करणे,
म्हणजे $\Setabs{x}{\phi(x)}$ या आकाराची व्याख्या करणे, पाहिले.
येथे आपण एखादा $\phi$ गुणधर्म वापरतो; या गुणधर्मात आधीच व्याख्या केलेल्या
संचांचा उल्लेख असू शकतो. उदाहरणार्थ, $A$ आणि $B$ हे संच असतील,
तर $\Setabs{x}{x \in A \lor x \in B}$ या संचात $A$ किंवा $B$ मधील
घटक असणाऱ्या सर्व वस्तू येतात. म्हणजेच, $A$ आणि $B$ मधील
घटक एकत्र करणारा हा संच आहे. \ref{sfr:set:uni:fig:union} मध्ये
दाखवल्याप्रमाणे त्याचे चित्र काढता येते; ठळक केलेला भाग $A$ आणि $B$
या दोन संचांतील घटक एकत्र दर्शवतो.

\begin{figure}
  \olasset{assets/diagrams/union.tikz}
  \caption{दोन संचांचा संयोग $A \cup B$ म्हणजे $A$ मधील आणि $B$ मधील
    घटक एकत्र घेऊन बनणारा संच.}
  \label{sfr:set:uni:fig:union}
\end{figure}

संच एकत्र करणारी ही क्रिया अतिशय उपयुक्त आहे आणि वारंवार वापरली जाते.
म्हणून तिला आपण आकारिक नाव आणि चिन्ह देतो.
\end{explain}

\begin{defn}[संयोग]
$A$ आणि $B$ या दोन संचांचा \emph{संयोग} $A \cup B$ असा लिहितात.
तो $A$, $B$ किंवा दोन्ही संचांतील घटक असणाऱ्या सर्व वस्तूंचा संच आहे.
\[A \cup B = \Setabs{x}{x \in A \lor x \in B}
\]
\end{defn}

\begin{ex}
तेच घटक किती वेळा आले आहेत याला महत्त्व नसल्याने, दोन संचांत
सामाईक घटक असेल, तर त्यांच्या संयोगात तो घटक एकदाच येतो.
उदाहरणार्थ, $\{ a, b, c\} \cup \{ a, 0, 1\} = \{a, b, c, 0, 1\}$.

संच आणि त्याचा एखादा उपसंच यांचा संयोग म्हणजे मोठा संचच:
$\{a,
b, c \} \cup \{a \} = \{a, b, c\}$.

संचाचा रिक्त संचाशी संयोग हा मूळ संचाशी समान असतो:
$\{a,
b, c \} \cup \emptyset = \{a, b, c \}$.
\end{ex}

\begin{prob}
$A \subseteq B$ असेल, तर $A \cup B = B$ हे सिद्ध करा.
\end{prob}

\begin{explain}
संयोगाच्या ``द्वैत'' क्रियेचाही विचार करता येतो. या क्रियेत असे सर्व
घटक घेतले जातात की जे $A$ मधील घटक आणि
$B$ मधीलही घटक आहेत, आणि त्यांचा संच बनवला जातो.
या क्रियेला \emph{छेद} म्हणतात. तिचे चित्र \ref{sfr:set:uni:fig:intersection} प्रमाणे काढता येते.
\begin{figure}
  \olasset{assets/diagrams/intersection.tikz}
  \caption{दोन संचांचा छेद $A \cap B$ म्हणजे त्यांतील सामाईक
    घटक यांचा संच.}
  \label{sfr:set:uni:fig:intersection}
\end{figure}
\end{explain}

\begin{defn}[छेद]
$A$ आणि $B$ या दोन संचांचा \emph{छेद} $A \cap B$ असा लिहितात.
तो $A$ आणि $B$ या दोन्ही संचांतील घटक असणाऱ्या सर्व वस्तूंचा संच आहे.
\[A \cap B = \Setabs{x}{x \in A \land x \in B}
\]
दोन संचांचा छेद रिक्त असेल, तर त्यांना \emph{विभक्त} म्हणतात.
याचा अर्थ, त्यांच्यात सामाईक घटक नसतात.
\end{defn}

\begin{ex}
दोन संचांत सामाईक घटक नसतील, तर त्यांचा छेद रिक्त असतो:
$\{ a, b, c\} \cap \{ 0, 1\} = \emptyset$.

दोन संचांत सामाईक घटक असतील, तर त्यांचा छेद म्हणजे त्या सर्वांचा संच:
$\{a, b, c \} \cap \{a, b, d \} = \{a, b\}$.

संच आणि त्याचा एखादा उपसंच यांचा छेद म्हणजे लहान संचच:
$\{a, b, c\} \cap \{a, b\} = \{a, b\}$.

कोणत्याही संचाचा रिक्त संचाशी छेद रिक्त असतो:
$\{a, b, c \}
\cap \emptyset = \emptyset$.
\end{ex}

\begin{prob}
$A \subseteq B$ असेल, तर $A \cap B = A$ हे काटेकोरपणे सिद्ध करा.
\end{prob}

\begin{explain}
दोनपेक्षा अधिक संचांचा संयोग किंवा छेदही बनवता येतो.
हे सर्वसाधारणपणे हाताळण्याची एक सुटसुटीत पद्धत पुढीलप्रमाणे आहे:
ज्या सर्व संचांचा संयोग (किंवा छेद) करायचा आहे, ते सर्व एका संचात एकत्र केले आहेत असे समजा.
मग या संचाच्या किमान एका घटक मध्ये असणाऱ्या सर्व वस्तूंचा संच
म्हणजे मूळ सर्व संचांचा संयोग अशी व्याख्या करता येते.
तसेच, या संचाच्या प्रत्येक घटक मध्ये असणाऱ्या सर्व वस्तूंचा संच
म्हणजे त्यांचा छेद अशी व्याख्या करता येते.
\end{explain}

\begin{defn}
$A$ हा संचांचा संच असेल, तर $\bigcup A$ म्हणजे $A$ मधील
घटक यांत असणाऱ्या घटक यांचा संच:
\begin{align*}
\bigcup A & = \Setabs{x}{x \text{ ही पुढील संचाच्या एका घटकात असते: } A},
\text{ म्हणजे,}\\
& = \Setabs{x}{\text{असा } B \in A
  \text{ असतो की } x \in B}
\end{align*}
\end{defn}

\begin{defn}
$A$ हा संचांचा संच असेल, तर $\bigcap A$ म्हणजे $A$ मधील सर्व घटकांत
सामाईक असणाऱ्या वस्तूंचा संच:
\begin{align*}
\bigcap A & = \Setabs{x}{x \text{ ही पुढील संचाच्या प्रत्येक घटकात असते: } A},
\text{ म्हणजे,}\\
 & = \Setabs{x}{\text{प्रत्येक } B \in A, x \in B}
\end{align*}
\end{defn}

\begin{ex}
$A = \{ \{ a, b \}, \{ a, d, e \}, \{ a, d \} \}$ असे समजा.
मग $\bigcup A = \{ a, b, d, e \}$ आणि $\bigcap A = \{ a \}$.
\end{ex}
\begin{prob}
	$A$ हा संच असेल आणि $A \in B$ असेल, तर $A \subseteq \bigcup B$ हे दाखवा.
\end{prob}

$A_1$, $A_2$, \dots या संचांच्या अनुक्रमासाठीही हेच करता येईल:
\begin{align*}
\bigcup_i A_i & = \Setabs{x}{x \text{ ही पुढीलपैकी एखाद्या संचात असते: } A_i}\\
\bigcap_i A_i & = \Setabs{x}{x \text{ ही पुढील प्रत्येक संचात असते: } A_i}.
\end{align*}

संचांसाठी \emph{निर्देशांक संच} दिलेला असेल, म्हणजे एखादा $I$ संच असा असेल की,
प्रत्येक $i \in I$ साठी आपण $A_i$ चा विचार करत असू, तर पुढील संक्षिप्त
रूपेही वापरता येतात:
\begin{align*}
	\bigcup_{i \in I} A_i & = \bigcup \Setabs{A_i }{i \in I}\\
	\bigcap_{i \in I} A_i & = \bigcap\Setabs{A_i}{i \in I}
\end{align*}

शेवटी, $A$ मध्ये असणारे पण $B$ मध्ये नसणारे सर्व घटक यांच्या
संचाचा विचार करायचा असू शकतो. त्याचे चित्र \ref{sfr:set:uni:difference} प्रमाणे काढता येते.

\begin{figure}
  \olasset{assets/diagrams/difference.tikz}
  \caption{दोन संचांचा फरक $A \setminus B$ म्हणजे $A$ मधील असे
    घटक की जे $B$ मधील घटक नाहीत, यांचा संच.}
  \label{sfr:set:uni:difference}
\end{figure}

\begin{defn}[फरक]
\emph{संचांचा फरक} $A \setminus B$ म्हणजे $A$ मधील असे सर्व
घटक की जे $B$ मधील घटक नाहीत, यांचा संच. म्हणजेच,
\[A\setminus B = \Setabs{x}{x\in A \text{ आणि } x \notin B}.
\]
\end{defn}

\begin{prob}
	$A \subsetneq B$ असेल, तर $B \setminus A \neq \emptyset$ हे सिद्ध करा.
\end{prob}











\section{जोड्या, क्रमित घटकसमूह आणि कार्तीय गुणाकार}\label{sfr:set:pai:sec}

\begin{explain}
संचांतील घटकांना क्रम नसतो, हे विस्तारात्मकतेवरून निष्पन्न होते.
म्हणून क्रम दर्शवायचा असेल, तर आपण \emph{क्रमित जोड्या} $\tuple{x, y}$ वापरतो.
अक्रमित जोडी $\{x, y\}$ मध्ये क्रमाला महत्त्व नसते:
$\{x, y\} = \{y, x\}$. क्रमित जोडीमध्ये मात्र क्रम महत्त्वाचा असतो:
$x
\neq y$ असेल, तर $\tuple{x, y} \neq \tuple{y, x}$.

संच सिद्धांतात क्रमित जोड्यांचा विचार कसा करावा? दोन क्रमित जोड्यांचे
पहिले घटक समान असतील आणि दुसरे घटकही समान असतील तेव्हा आणि केवळ तेव्हाच
त्या जोड्या समान असतात, ही कल्पना जपणे आवश्यक आहे. म्हणजेच:
\[  \tuple{a, b}= \tuple{c, d}\text{ तेव्हा आणि केवळ तेव्हाच, जेव्हा }a = c \text{ आणि }b=d.
\]
वीनर--कुराटोव्स्की यांच्या व्याख्येने संच सिद्धांतात क्रमित जोड्यांची व्याख्या करता येते.
\end{explain}

\begin{defn}[क्रमित जोडी]\label{sfr:set:pai:wienerkuratowski}
	$\tuple{a, b} = \{\{a\}, \{a, b\}\}$.
\end{defn}

\begin{prob}
	\ref{sfr:set:pai:wienerkuratowski} वापरून, $\tuple{a,
	b}= \tuple{c, d}$ तेव्हा आणि केवळ तेव्हाच, जेव्हा $a = c$ आणि $b=d$
दोन्ही असतात, हे सिद्ध करा.
\end{prob}

\begin{explain}
क्रमित जोडीची व्याख्या निश्चित केल्यावर तिच्या साहाय्याने आणखी संचांची
व्याख्या करता येते. उदाहरणार्थ, कधीकधी दोनपेक्षा अधिक वस्तूंचे क्रमित
अनुक्रम हवे असतात: \emph{त्रिके} $\tuple{x, y, z}$,
\emph{चतुष्के} $\tuple{x, y, z, u}$ इत्यादी.
त्रिक म्हणजे ज्याचा पहिला घटक स्वतःच क्रमित जोडी आहे अशी विशिष्ट क्रमित
जोडी असे मानता येते: $\tuple{x, y, z}$ म्हणजे $\tuple{\tuple{x, y},z}$.
चतुष्कांसाठीही हेच लागू होते: $\tuple{x,y,z,u}$ म्हणजे
$\tuple{\tuple{\tuple{x,y},z},u}$, आणि असेच पुढे.
सर्वसाधारणपणे आपण \emph{क्रमित $n$-घटकसमूह} $\tuple{x_1, \dots, x_n}$ यांचा उल्लेख करतो.

क्रमित जोड्यांचे, किंवा इतर क्रमित $n$-घटकसमूहांचे, काही संच उपयुक्त ठरतील.
\end{explain}

\begin{defn}[कार्तीय गुणाकार]
$A$ आणि $B$ हे संच दिलेले असतील, तर त्यांचा \emph{कार्तीय गुणाकार}
$A \times B$ याची व्याख्या अशी करतात:
\[  A \times B = \Setabs{\tuple{x, y}}{x \in A \text{ आणि } y \in B}.
\]
\end{defn}

\begin{ex}
$A = \{0, 1\}$ आणि $B = \{1, a, b\}$ असतील, तर त्यांचा गुणाकार असा आहे:
\[A \times B = \{ \tuple{0, 1}, \tuple{0, a}, \tuple{0, b},
    \tuple{1, 1}, \tuple{1, a}, \tuple{1, b} \}.
\]
\end{ex}

\begin{ex}
$A$ हा संच असेल, तर $A$ चा स्वतःशी गुणाकार $A \times A$ हा $A^2$
असाही लिहितात. तो $x, y \in A$ असणाऱ्या \emph{सर्व} $\tuple{x, y}$ जोड्यांचा
संच आहे. सर्व $\tuple{x, y, z}$ त्रिकांचा संच $A^3$ असतो, आणि असेच पुढे.
पुढील पुनरावर्ती व्याख्या देता येते:
\begin{align*}
  A^1 & = A\\
  A^{k+1} & = A^k \times A
\end{align*}
\end{ex}

\begin{prob}
$\{1, 2, 3\}^3$ मधील सर्व घटक लिहा.
\end{prob}

\begin{prop}\label{sfr:set:pai:cardnmprod}
$A$ मध्ये $n$ घटक आणि $B$ मध्ये $m$ घटक असतील,
तर $A
\times B$ मध्ये $n\cdot m$ घटक असतात.
\end{prop}

\begin{proof}
$A$ मधील प्रत्येक घटक $x$ साठी, $\tuple{x, y} \in A \times B$
या आकाराचे $m$ घटक असतात.
$B_x = \Setabs{\tuple{x, y}}{y
  \in B}$ असे मानू. $x_1 \neq x_2$ असेल तेव्हा $\tuple{x_1, y} \neq
\tuple{x_2, y}$ असल्यामुळे $B_{x_1} \cap B_{x_2} = \emptyset$.
पण $A = \{x_1,
\dots, x_n\}$ असेल, तर $A \times B = B_{x_1} \cup \dots \cup B_{x_n}$;
म्हणून त्यात $n\cdot m$ घटक असतात.

हे चित्ररूपाने पाहण्यासाठी $A \times B$ मधील घटक जाळीच्या स्वरूपात मांडा:
\[\begin{array}{rcccc}
  B_{x_1} = & \{\tuple{x_1, y_1} & \tuple{x_1, y_2} & \dots & \tuple{x_1, y_m}\}\\
  B_{x_2} = & \{\tuple{x_2, y_1} & \tuple{x_2, y_2} & \dots & \tuple{x_2, y_m}\}\\
  \vdots & & \vdots\\
  B_{x_n} = & \{\tuple{x_n, y_1} & \tuple{x_n, y_2} & \dots & \tuple{x_n, y_m}\}
\end{array}
\]
सर्व $x_i$ वेगवेगळे आहेत आणि सर्व $y_j$ वेगवेगळे आहेत. त्यामुळे या जाळीतील
कोणत्याही दोन जोड्या समान नाहीत, आणि त्यांची संख्या $n\cdot m$ आहे.
\end{proof}

\begin{prob}
$k$ वरील गणिती विगमनाने असे दाखवा की, प्रत्येक $k \ge 1$ साठी,
$A$ मध्ये $n$ घटक असतील, तर $A^k$ मध्ये $n^k$ घटक असतात.
\end{prob}

\begin{ex}
$A$ हा संच असेल, तर $A$ वरील \emph{शब्द} म्हणजे $A$ मधील
घटक यांचा कोणताही अनुक्रम.
असा अनुक्रम म्हणजे $A$ मधील घटक यांचा $n$-घटकसमूह असे मानता येते.
उदाहरणार्थ, $A = \{a, b, c\}$ असेल, तर ``$bac$'' या अनुक्रमाचा
$\tuple{b, a, c}$ हे त्रिक म्हणून विचार करता येतो.
शब्दांना, म्हणजे चिन्हांच्या अनुक्रमांना, संगणकशास्त्रात अत्यंत महत्त्व आहे.
संकेतानुसार, $A$ मधील घटक यांना आपण $1$ लांबीचे अनुक्रम मानतो,
आणि $\emptyset$ याला $0$ लांबीचा अनुक्रम मानतो.
मग $A$ वरील \emph{सर्व} शब्दांचा संच असा आहे:
\[A^* = \{\emptyset\} \cup A \cup A^2 \cup A^3 \cup \dots
\]
\end{ex}












\section{रसेलची विरोधापत्ती}\label{sfr:set:rus:sec}

विस्तारात्मकतेमुळे $\Setabs{x}{\phi(x)}$ ही चिन्हपद्धती $\phi(x)$ असणाऱ्या
सर्व $x$ च्या \emph{त्या विशिष्ट} संचासाठी वापरणे योग्य ठरते.
मात्र, विस्तारात्मकतेमुळे \emph{प्रत्यक्षात} फक्त पुढील विचार योग्य ठरतो.
ज्याचे सदस्य सर्व आणि केवळ $\phi$ असणाऱ्या वस्तू आहेत असा संच \emph{असेल},
\emph{तर} असा एकच संच असतो. दुसऱ्या शब्दांत: एखादा $\phi$ निश्चित केल्यावर,
$\Setabs{x}{\phi(x)}$ हा संच \emph{अस्तित्वात असल्यास} तो एकमेव असतो.

पण ही अट महत्त्वाची आहे!{} प्रत्येक गुणधर्मापासून \emph{गुणधर्माधारित संचरचना}
करता येतेच असे नाही. म्हणजेच, काही गुणधर्म संचांची व्याख्या करत \emph{नाहीत}.
सर्वच गुणधर्मांनी संचांची व्याख्या होत असती, तर आपल्याला थेट व्याघातांना
सामोरे जावे लागले असते. याचे सर्वांत प्रसिद्ध उदाहरण म्हणजे रसेलची विरोधापत्ती.

संच हे इतर संचांचे घटक असू शकतात; उदाहरणार्थ, $A$ या संचाचा घातसंच
हा संचांनी बनलेला असतो. म्हणून एखादा संच दुसऱ्या संचाचा घटक आहे का,
असा प्रश्न विचारणे किंवा त्याचा शोध घेणे अर्थपूर्ण आहे. संच स्वतःचाच सदस्य
असू शकतो का? संचाच्या कल्पनेत असे काही दिसत नाही की ज्यामुळे ही शक्यता
नाकारली जाईल. उदाहरणार्थ, \emph{सर्व} संच मिळून वस्तूंचा समूह बनत असेल,
तर त्या सर्वांना एका संचात एकत्र करता येईल असे वाटू शकते; तो म्हणजे सर्व
संचांचा संच. तो स्वतः एक संच असल्यामुळे सर्व संचांच्या संचाचा घटक असेल.

स्वतःला घटक म्हणून न सामावणाऱ्या, म्हणजे \emph{स्वतःचे सदस्य नसणाऱ्या}
संचांचा गुणधर्म विचारात घेतल्यावर रसेलची विरोधापत्ती उद्भवते.
स्वतःला घटक म्हणून न सामावणाऱ्या सर्व संचांचा एक संच आहे असे आपण
मानले, तर काय होईल? असा
\[R = \Setabs{x}{x \notin x}
\]
संच अस्तित्वात आहे का? तो अस्तित्वात नसतो हे आपण सिद्ध करू शकतो.

\begin{thm}[रसेलची विरोधापत्ती]\label{sfr:set:rus:thm:russells-paradox}
	$R = \Setabs{x}{x \notin x}$ असा कोणताही संच नाही.
\end{thm}

\begin{proof}
$R = \Setabs{x}{x \notin x}$ अस्तित्वात असेल, तर
$R \in R$ तेव्हा आणि केवळ तेव्हाच, जेव्हा $R \notin R$; हा व्याघात आहे.
\end{proof}


\begin{explain}
ही सिद्धता अधिक सावकाश पाहू. $R$ अस्तित्वात असेल, तर $R \in
R$ आहे की नाही हा प्रश्न अर्थपूर्ण आहे. खरोखर $R \in R$ आहे असे समजा.
आता, स्वतःचे घटक नसणाऱ्या सर्व संचांचा संच अशी $R$ ची व्याख्या केली होती.
त्यामुळे $R \in R$ असेल, तर $R$ मध्ये स्वतःच $R$ चा व्याख्यक गुणधर्म नसतो.
पण हा गुणधर्म असणारेच संच $R$ मध्ये असतात. म्हणून $R$ हा $R$ चा
घटक असू शकत नाही; म्हणजेच $R \notin R$.
परंतु $R$ हा $R$ चा घटक आहे आणि नाही, असे दोन्ही होऊ शकत नाही;
म्हणून व्याघात मिळतो.

$R \in R$ या गृहीतकामुळे व्याघात येत असल्याने $R \notin R$ मिळते.
पण यामुळेही व्याघात येतो!{} कारण $R
\notin R$ असेल, तर $R$ मध्ये स्वतःच $R$ चा व्याख्यक गुणधर्म असतो;
म्हणून स्वतःचे सदस्य नसणाऱ्या इतर सर्व संचांप्रमाणे $R$ हा $R$ चा
घटक असेल. आणि पुन्हा, $R$ चा घटक नसणे आणि असणे
या दोन्ही गोष्टी एकत्र शक्य नाहीत.
\end{explain}


\begin{digress}
रसेलची विरोधापत्ती टाळणारा, म्हणजे $R = \Setabs{x}{x \notin x}$
अस्तित्वात आहे हा \emph{विसंगत} दावा न करणारा संच सिद्धांत कसा मांडायचा?
त्यासाठी संच केव्हा अस्तित्वात असतात (आणि केव्हा नसतात) हे सांगणाऱ्या
अगदी नेमक्या अटी देणारी स्वयंसिद्धके मांडावी लागतील.

या प्रकरणात आराखडा दिलेला संच सिद्धांत असे करत नाही. तो
\emph{खरोखरच अनौपचारिक} आहे. त्यातून एवढेच सांगितले जाते की संच
विस्तारात्मकतेचे पालन करतात आणि काही संच दिलेले असतील, तर त्यांचा संयोग,
छेद इत्यादी बनवता येतात. संच सिद्धांत यापेक्षा अधिक काटेकोरपणे विकसित करणे
शक्य आहे.
\end{digress}




